We design three guidance heuristics ($h_1$, $h_2$, $h_3$) to steer A* toward executable arbitrage routes. Each captures a different dimension of execution feasibility via an additive penalty $h(n) \geq 0$ in $f(n) = g(n) + h(n)$. These are \emph{guidance penalties}, not admissible lower bounds. We also evaluate a parallelized multi-start baseline that wraps any heuristic.

\subsection{$h_1$ -- Liquidity Heuristic}

The liquidity heuristic penalizes low-volume markets that cannot support the required trade size. We compute a liquidity ratio from 24-hour quote volume, the remaining arbitrage window $T_{\textit{remain}}$, and order size:
\[
\textit{liquidity\_ratio} = \frac{(\textit{24h\_vol} / 86400) \cdot T_{\textit{remain}}}{\textit{order\_size\_usd}},
\quad
h_1(n) = \lambda_{liq}\!\left(1 - \textit{clamp}\!\left(\tfrac{\log_{10}(\textit{liquidity\_ratio})+1}{2},\,0,\,1\right)\right).
\]
This heuristic favours deep, actively traded markets and is time-aware.

\subsection{$h_2$ -- Slippage Heuristic}

The slippage heuristic models execution price impact using real-time order-book depth. Given mid-price $P_{mid}$ and the VWAP required to fill order size $Q$:
\[
h_2(n) = \lambda_{slip} \cdot \max\!\bigl\{0,\; \textit{slippage}_{bps} - \theta\bigr\},
\quad
\textit{slippage}_{bps} = \frac{\textit{VWAP}(Q) - P_{mid}}{P_{mid}} \times 10^4,
\]
where $\theta$ is a configurable tolerance. Edge costs encode static taker and withdrawal fees; $h_2$ captures \emph{dynamic} order-book slippage, so there is no double-counting with $g(n)$.

\subsection{$h_3$ -- Chain Congestion and Exchange Reliability}

$h_3$ penalizes transfer latency and operational risk using pre-computed chain transfer times and static reliability scores $s(n) \in [0,1]$:
\[
h_3(n) = \lambda_{chain}\!\left(\frac{t_{\min}}{T_{\text{remain}}}\right) + \lambda_{exchange}\!\left(1 - s(n)\right).
\]

\subsection{Parallelized Multi-Start Baseline}

The parallelized baseline runs $k$ independent A* searches in parallel from randomly sampled starting nodes, each using a base heuristic, and returns the most profitable result.
